由维数可知,
\[
\pi_*(\xi^j)=0\qquad\text{若 }j<a-1,
\]
而
\[
\pi_*(\xi^{a-1})=1_X
\]
这是因为几何上 $\xi^{a-1}$ 与每个纤维相交于一个点. 因此投影公式给出
\begin{equation}\tag{5.3.9}
\pi_*\!\left(
\sum_{i=0}^{a-1}\pi^*(U_i)\xi^i
\right)=U_{a-1}.
\end{equation}
\footnote{印本所列公式混用了下标 $i,j$, 并将右边写为 $U_{j-i}$. 紧接在前面的两个推出恒等式确定了系数应为 $U_{a-1}$.}
由 (5.3.9), 鉴于 (5.3.1), 用 $T$ 替换 $\xi$, 用 $S_i$ 替换 $C_i$, 则断言 (5.3.5) 由下一引理得出.

\subsubsection*{引理 5.4.}
\underline{设 $a>0$ 为整数, $A=\mathbb Z[[X_1,\ldots,X_a]]$ 为 $a$ 个不定元的形式幂级数环. 设 $B$ 为 $A$-代数}
\[
B=A[[T]]\bigg/\left(\prod_{i=1}^{a}(T-X_i)\right).
\]
\underline{则 $B$ 为自由 $A$-模, 基为 $1,T,\ldots,T^{a-1}$:}
\begin{equation}\tag{5.4.1}
B=AT^{a-1}\oplus AT^{a-2}\oplus\cdots\oplus A.
\end{equation}
\underline{设 $S_i\in A$ 为第 $i$ 个初等对称函数}
\[
S_i=\sum_{1\leq\alpha_1<\cdots<\alpha_i\leq a}
X_{\alpha_1}\cdots X_{\alpha_i},
\]
\underline{从而}
\begin{equation}\tag{5.4.2}
\sum_{i=0}^{a}(-1)^iS_iT^{a-i}=0
\qquad\text{在 }B.
\end{equation}
\underline{元素 $1+S_1+\cdots+S_a$ 在 $A$ 中可逆; 将其逆写为}
\begin{equation}\tag{5.4.3}
\frac{1}{\sum_{i=0}^{a}S_i}
=\sum_{j\geq0}b_j,
\qquad
b_j\text{ 为齐次元素, 次数为 }j\text{, 位于 }A,
\qquad b_0=1.
\end{equation}
\underline{对每个整数 $j\geq0$, 利用 (5.4.1) 写成}
\begin{equation}\tag{5.4.4}
T^{a-1+j}
=a_jT^{a-1}+{}?\,T^{a-2}+{}?\,T^{a-3}+\cdots+{}?,
\qquad ?\in A.
\end{equation}
\underline{由此定义 $j$ 次齐次元素 $a_j\in A$. 则对每个 $j\geq0$,}
\begin{equation}\tag{5.4.5}
a_j=(-1)^jb_j.
\end{equation}

